% Generated draft: requires mathematical and source review.
\documentclass[11pt]{article}
\usepackage[margin=25mm]{geometry}
\usepackage{fontspec}
\setmainfont{Latin Modern Roman}
\usepackage{amsmath,amssymb,mathtools,unicode-math}
\setmathfont{Latin Modern Math}
\usepackage{enumitem,needspace,xurl,hyperref}
\setlength{\parindent}{0pt}
\setlength{\parskip}{5pt}
\setlength{\emergencystretch}{3em}
\newcommand{\N}{\mathbb{N}}
\newcommand{\Z}{\mathbb{Z}}
\newcommand{\Q}{\mathbb{Q}}
\newcommand{\R}{\mathbb{R}}
\newcommand{\C}{\mathbb{C}}
\newcommand{\F}{\mathbb{F}}
\newcommand{\A}{\mathbb{A}}
\newcommand{\PP}{\mathbb{P}}
\newcommand{\E}{\mathbb{E}}
\newcommand{\eps}{\varepsilon}
\newcommand{\dd}{\,\mathrm{d}}
\begin{document}
\section*{UnsolvedMath ID: 6800018}
\textbf{Source problem code:} catalogue entry
\subsection{Definitions and mathematical statement}
Equip the complex projective plane $M=\C P^2$ with the Fubini--Study metric $g_{FS}$ identified by the research summary. This is a compact real four-manifold; the scale normalization is not stated. Write $V=\operatorname{vol}_{g_{FS}}(M)$. For $0<v<V$, define the isoperimetric profile
\[
 I(v)=\inf\{P_{g_{FS}}(E):E\subseteq M\text{ of finite perimeter},
                         \ \operatorname{vol}_{g_{FS}}(E)=v\},
\]
where $P_{g_{FS}}(E)=\mathcal H^3_{g_{FS}}(\partial^*E)$ in the finite-perimeter convention.

A geodesic sphere is a distance sphere $\{x:d_{g_{FS}}(p,x)=r\}$ at a radius with the intended regularity. The conjecture asks whether, for every prescribed $v\in(0,V)$, a region bounded by such a sphere attains $I(v)$. In a ball/complement interpretation, this means finding a center $p$, radius $r$, and the appropriate side $E$ of that sphere with
\[
 \operatorname{vol}_{g_{FS}}(E)=v,\qquad
 P_{g_{FS}}(E)=I(v).
\]
The side convention should be confirmed: on a compact manifold the complement has the same boundary and complementary volume. The assertion supplies a spherical minimizing candidate at each volume; it does not explicitly claim uniqueness or that every minimizer is spherical.

The original statement does not identify the metric, normalization, admissible regularity, or treatment of cut-locus radii. The Fubini--Study and finite-perimeter formulations are explicit standard interpretations for review, not hidden assumptions extending the problem to arbitrary metrics on $\C P^2$.
\subsection{Short English statement}
Do geodesic spheres solve the prescribed-volume least-perimeter problem in the Fubini--Study complex projective plane? The question asks for global minimality at every nontrivial volume, with no extra uniqueness claim.

The supplied summary describes geodesic spheres as candidates and mentions stability and small/large-volume issues, but provides no radius range, exact theorem, or citation for stability. It says a complete all-volume result was not verified. The partial-progress label and 2018 question are retained with current status marked for review.
\subsection{Sources}
\begin{itemize}
\item Local UnsolvedMath registry, permanent ID 6800018, source code catalogue entry, ``Isoperimetric Problem in $\C P^2$,'' proposed by Frank Morgan, Williams College, Williamstown, 2018. Dataset curation and metadata: Ulam AI and UnsolvedMath Contributors, CC BY 4.0.
\item Morgan and Pansu, \emph{A List of Open Problems in Differential Geometry} (2018), Question environment 18, \url{https://www.imo.universite-paris-saclay.fr/~pansu/problems_MTDG.tex}; source-TeX extraction, accessed 2026-07-29.
\item The summary mentions Morgan's analysis, geodesic-sphere stability, complex-space-form results, and Hopf-fibration context without complete references or exact hypotheses.
\item Source-specific provenance: the item is labelled a question, current status NEEDS\_REVIEW; the public source TeX is on a coauthor's publication page, with redistribution terms NEEDS\_REVIEW, separately from the dataset licence.
\end{itemize}
\end{document}
